\section{Verificación Formal de Protocolos de Seguridad (Modelo Dolev--Yao, Tamarin Prover)}
% =============================================================================

\subsection{Fundamentación de la Verificación Formal en Protocolos Criptográficos}

Un protocolo criptográfico es el conjunto ordenado de reglas y pasos que definen cómo se intercambian mensajes dos o más computadoras a través de una red para llevar a cabo una tarea segura, como iniciar sesión, realizar un pago o emitir un voto.

En la práctica de la ciberseguridad, se ha comprobado repetidamente que un sistema puede utilizar los algoritmos criptográficos más fuertes del mundo (como ML-KEM o ML-DSA) y, sin embargo, ser inseguro si el orden o la estructura de los mensajes del protocolo tiene un error de diseño (Cremers y Mauw, 2012). Por ejemplo, si una puerta blindada de acero (el algoritmo criptográfico) se instala en un marco mal construido que permite sacar las bisagras desde afuera (el protocolo), un atacante no necesitará romper el acero para entrar.

Un caso clásico en la historia de la Informática es el protocolo de autenticación de Needham-Schroeder, el cual fue considerado seguro durante casi diecisiete años hasta que se descubrió, mediante una herramienta de análisis formal, que un atacante podía interceptar y reenviar mensajes válidos en distinto orden para hacerse pasar por otro usuario. Asimismo, en el área de votación electrónica, Baloglu et al. (2021) utilizaron verificación formal para encontrar fallas de duplicación de boletas en los sistemas Helios y Belenios que no habían sido detectadas mediante revisiones manuales.

Cuando un programador realiza pruebas de software tradicionales (como pruebas unitarias o de integración), únicamente comprueba que el programa funciona bien en los casos que él mismo imaginó, pero nunca puede probar todos los ataques posibles. En cambio, la \textbf{verificación formal} consiste en representar el protocolo mediante lógica matemática para que un motor automatizado explore todas las combinaciones posibles de mensajes y demuestre con certeza que ningún ataque puede vulnerar las reglas definidas (Schmidt et al., 2012).

\subsection{El Modelo de Atacante Dolev-Yao y el Conocimiento Adversarial}

Para comprobar la seguridad de un protocolo mediante verificación formal, primero es necesario definir qué capacidades tiene el atacante. En ciberseguridad, el estándar más exigente y utilizado a nivel mundial es el \textbf{modelo de atacante Dolev-Yao}, formulado por Danny Dolev y Andrew Yao (1983).

La idea principal del modelo Dolev-Yao es asumir el peor escenario posible: \textit{el atacante tiene control total sobre toda la red de telecomunicaciones}. Es decir, cualquier mensaje que viaje entre el navegador del votante y el servidor electoral pasa por las manos del atacante, quien cuenta con las siguientes capacidades y límites:

\begin{itemize}
	\item \textbf{Control completo de los mensajes en la red:} El atacante puede leer todos los paquetes enviados, guardar copias, bloquear mensajes para que no lleguen a su destino, modificarlos, cambiarles el orden o volver a enviar mensajes antiguos mediante ataques de repetición.

	\item \textbf{Creación de nuevos mensajes:} El atacante puede fabricar mensajes falsos combinando cualquier dato que haya capturado previamente en la red con datos inventados por él mismo.

	\item \textbf{Respeto a la criptografía perfecta:} La única limitación del atacante Dolev-Yao es que no puede adivinar claves por magia ni romper las matemáticas de los algoritmos: no puede descifrar un mensaje cifrado si no posee la clave privada correspondiente, no puede falsificar una firma digital sin tener la clave secreta de firma, y no puede invertir un resumen hash.
\end{itemize}

Dentro de la herramienta Tamarin Prover, todo lo que el atacante va aprendiendo se representa mediante el predicado especial:
\begin{equation}
	K(m)
\end{equation}
el cual significa que ``el atacante conoce el dato $m$''. Cada vez que un participante envía un mensaje a la red mediante la instrucción de salida, ese dato $t$ pasa inmediatamente al conocimiento del atacante ($K(t)$). De igual manera, cualquier dato recibido de la red mediante la instrucción de entrada puede haber sido alterado o fabricado por el atacante usando todo lo que conoce hasta ese momento.

Aplicar el modelo Dolev-Yao a la votación electrónica universitaria significa demostrar que, aunque un intruso controle por completo la red WiFi del lugar o los enrutadores de Internet, no podrá leer el contenido de ningún voto ni introducir boletas falsas en la urna.

\subsection{Criptografía Simbólica versus Criptografía Computacional}

En la investigación en seguridad informática existen dos maneras complementarias de analizar y demostrar que un protocolo es seguro: el enfoque computacional y el enfoque simbólico.

\begin{enumerate}
	\item \textbf{Enfoque Computacional:} Estudia los algoritmos trabajando directamente con cadenas de unos y ceros ($\{0, 1\}^*$) y probabilidades matemáticas. Demuestra que a un atacante le tomaría miles de años de cálculo resolver problemas difíciles como Module-LWE o Module-SIS. Es el enfoque que respalda la seguridad interna de los algoritmos ML-KEM y ML-DSA.

	\item \textbf{Enfoque Simbólico:} Trata a los algoritmos criptográficos como ``cajas cerradas ideales'' regidas por ecuaciones exactas y analiza el flujo lógico de los mensajes entre los participantes para descubrir si existe algún hueco en las reglas del protocolo (Basin et al., 2015). Es el enfoque que utiliza la herramienta Tamarin Prover.
\end{enumerate}

En la Tabla \ref{tab:simbolico_vs_computacional} se resumen las diferencias principales entre ambos enfoques de análisis de seguridad.

\begin{table}[htbp]
	\centering
	\small
	\caption{Comparación metodológica entre criptografía simbólica y computacional.}
	\label{tab:simbolico_vs_computacional}
	\begin{tabular}{p{2.8cm} p{5.3cm} p{5.5cm}}
		\toprule
		\textbf{Dimensión} & \textbf{Enfoque Simbólico (Tamarin)} & \textbf{Enfoque Computacional} \\
		\midrule
		\textbf{Representación} & Símbolos y términos algebraicos & Cadenas binarias de bits $\{0, 1\}^*$ \\
		\textbf{Primitivas} & Cajas negras gobernadas por ecuaciones & Algoritmos probabilísticos en tiempo polinomial \\
		\textbf{Resultado} & Demostración exacta o traza de ataque & Probabilidad de éxito despreciable $\epsilon$ \\
		\textbf{Alcance} & Número ilimitado de sesiones en paralelo & Acotado por recursos de cómputo \\
		\textbf{Herramienta} & Verificación automatizada (Tamarin / Maude) & Reducciones matemáticas manuales \\
		\bottomrule
	\end{tabular}
	\par\vspace{3pt}{\footnotesize\textit{Fuente: Adaptado de Cremers y Mauw (2012).}}
\end{table}

Como muestra la Tabla \ref{tab:simbolico_vs_computacional}, ambos métodos se complementan en la presente tesis: los estándares del NIST garantizan la fortaleza computacional de los algoritmos, mientras que Tamarin Prover verifica simbólicamente que el protocolo electoral los utilice de manera segura. Para ello, las operaciones post-cuánticas se representan en Tamarin mediante ecuaciones algebraicas claras:

\begin{itemize}
	\item \textbf{Encapsulamiento de Claves (ML-KEM / FIPS 203):} Se representa con las funciones $\text{encaps}(pk, k)$ y $\text{decaps}(ct, sk)$, unidas por la regla:
	\begin{equation}
		\text{decaps}(\text{encaps}(pk(sk), k), sk) = k
	\end{equation}
	Esta ecuación indica que solo quien posee la clave privada $sk$ correspondiente a la clave pública $pk(sk)$ puede abrir la cápsula y recuperar la llave secreta $k$.

	\item \textbf{Firma Digital (ML-DSA / FIPS 204):} Se representa con las funciones $\text{sign}(m, sk)$ y $\text{verify}(\sigma, m, pk)$, unidas por la regla:
	\begin{equation}
		\text{verify}(\text{sign}(m, sk), m, pk(sk)) = \text{true}
	\end{equation}
	Esta ecuación expresa que la verificación devuelve verdadero únicamente cuando la firma fue generada con la clave privada auténtica $sk$ sobre el mismo mensaje $m$.

	\item \textbf{Cifrado Simétrico y Funciones Hash (AES-GCM y FIPS 202):} El cifrado simétrico cumple $\text{sdec}(\text{senc}(m, k), k) = m$ (se cifra y descifra con la misma llave $k$), mientras que la función hash $h(m)$ no tiene ecuación inversa, reflejando que es una operación unidireccional.
\end{itemize}

\subsection{Semántica Operacional y Arquitectura de Tamarin Prover}

Tamarin Prover (Schmidt et al., 2012; Meier et al., 2013) es una herramienta informática de código abierto desarrollada en la Escuela Politécnica Federal de Zúrich (ETH Zürich) para verificar automáticamente protocolos criptográficos. Su funcionamiento se basa en un modelo llamado \textbf{reescritura de multiconjuntos}.

De forma sencilla, se puede imaginar el estado del sistema en Tamarin como una ``mesa de trabajo'' donde se colocan tarjetas llamadas \textbf{hechos}. Cada hecho guarda información sobre lo que ocurre en un momento dado (por ejemplo, que un votante fue registrado o que una boleta fue guardada en la urna).

\subsubsection{Clasificación de Hechos en Tamarin}

En un modelo de Tamarin existen tres tipos de hechos que permiten representar con exactitud el comportamiento del sistema electoral:

\begin{enumerate}
	\item \textbf{Hechos Lineales (de un solo uso):} Son hechos que, una vez utilizados por un paso del protocolo, se consumen y desaparecen del estado. Esta propiedad resulta sumamente útil para la votación electrónica: por ejemplo, la autorización para votar de un ciudadano se modela como un hecho lineal; cuando el ciudadano emite su voto, ese hecho se consume y deja de existir, impidiendo que el mismo elector vuelva a votar por segunda vez.

	\item \textbf{Hechos Persistentes (permanentes):} Se escriben anteponiendo un signo de admiración. Representan datos públicos o registros permanentes que pueden consultarse muchas veces sin gastarse, tales como las claves públicas de las autoridades electorales o el padrón oficial.

	\item \textbf{Hechos Especiales del Sistema:} Tamarin incluye tres tipos de eventos propios para interactuar con el entorno:
	\begin{itemize}
		\item \textbf{Generación aleatoria fresca:} Genera un número o clave aleatoria completamente nueva que nadie ha visto antes.
		\item \textbf{Emisión de mensajes:} Envía el mensaje hacia la red pública.
		\item \textbf{Recepción de mensajes:} Recibe el mensaje desde la red pública.
	\end{itemize}
\end{enumerate}

\subsubsection{Estructura de las Reglas de Reescritura}

Cada paso del protocolo (como registrar a un votante, emitir un voto o contar una boleta) se programa en Tamarin mediante una \textbf{regla de reescritura} que tiene tres partes:
\begin{equation}
	\text{Regla}: \quad [l_1, \dots, l_n] \xrightarrow{[a_1, \dots, a_k]} [r_1, \dots, r_m]
\end{equation}
\begin{itemize}
	\item \textbf{Premisas ($[l_1, \dots, l_n]$):} Son las condiciones o hechos que deben existir previamente en el sistema para que el paso pueda ejecutarse.
	\item \textbf{Etiquetas de acción ($[a_1, \dots, a_k]$):} Son marcas de registro (como un historial o bitácora con marca de tiempo $\#t$) que indican qué acción acaba de ocurrir para poder revisarla después en las pruebas de seguridad.
	\item \textbf{Conclusiones ($[r_1, \dots, r_m]$):} Son los nuevos hechos que quedan creados en el sistema después de ejecutar la regla.
\end{itemize}

A la secuencia ordenada en el tiempo de todas las etiquetas de acción que ocurren durante una ejecución del protocolo se le denomina \textbf{traza de ejecución}:
\begin{equation}
	\text{Traza} = [A_1 @ \#t_1, \; A_2 @ \#t_2, \; \dots, \; A_N @ \#t_N] \quad \text{con} \quad \#t_1 < \#t_2 < \dots < \#t_N
\end{equation}

\subsection{Especificación Lógica: Restricciones y Lemas Formales}

Una vez escritas las reglas del protocolo, se debe indicar a Tamarin qué reglas de control debe respetar el sistema (restricciones) y qué propiedades de seguridad debe demostrar matemáticamente (lemas).

\subsubsection{Restricciones (Axiomas del Protocolo)}

Las restricciones representan las verificaciones obligatorias que hace el software (como comprobar que una firma digital sea válida o que un usuario no haya votado antes). Si una secuencia de pasos incumple una restricción, el servidor la rechaza de inmediato.

Por ejemplo, la restricción de unicidad del voto se escribe en lógica formal indicando que si aparece un voto registrado para el ciudadano $id$ en el instante de tiempo $\#i$, y también aparece un voto para el mismo ciudadano $id$ en el instante $\#j$, entonces ambos instantes deben ser obligatoriamente el mismo momento ($\#i = $\#j):
\begin{equation}
	\forall id, \#i, \#j. \; \big(\text{VotoUnico}(id) @ \#i \;\land\; \text{VotoUnico}(id) @ \#j\big) \implies \#i = \#j
\end{equation}
Es decir, el sistema prohíbe que existan dos momentos distintos en los que una misma persona deposite un voto.

\subsubsection{Lemas Formales: Sanidad y Seguridad}

Los lemas son los teoremas de seguridad que Tamarin Prover debe demostrar de forma automática. Se dividen en dos tipos:

\begin{enumerate}
	\item \textbf{Lemas de Ejecutabilidad:} Comprueban que existe al menos una traza normal donde los participantes honestos pueden completar todo el proceso electoral de principio a fin sin quedarse bloqueados.

	\item \textbf{Lemas de Seguridad Universal:} Comprueban que una propiedad de seguridad se cumple en \textit{absolutamente todas las ejecuciones posibles}, sin importar qué ataques intente realizar el adversario Dolev-Yao:
	\begin{itemize}
		\item \textbf{Confidencialidad del voto:} Establece que para todo voto $m$ emitido en un instante $\#t$, no existe ningún instante $\#t_2$ en el que el atacante llegue a conocer $m$:
		\begin{equation}
			\forall m, \#t. \; \text{SecretoVoto}(m) @ \#t \implies \neg (\exists \#t_2. \; K(m) @ \#t_2)
		\end{equation}

		\item \textbf{Integridad del escrutinio:} Establece que si un voto $m$ es contabilizado en el instante $\#t_c$, entonces forzosamente tuvo que ser emitido sin alteraciones en un instante anterior $\#t_e < \#t_c$:
		\begin{equation}
		\begin{split}
			\forall id, m, \#t_c. \; &\text{VotoContado}(id, m) @ \#t_c \implies \\
			&\big(\exists \#t_e < \#t_c. \; \text{VotoEmitido}(id, m) @ \#t_e\big)
		\end{split}
		\end{equation}

		\item \textbf{Autenticidad y no repetición o acuerdo inyectivo:} Garantiza que cada voto aceptado por la urna proviene de un elector legítimo acreditado y no es una copia repetida por un atacante en la red.
	\end{itemize}
\end{enumerate}

Si Tamarin comprueba que todos los casos cumplen la fórmula, emite el resultado verificado. Por el contrario, si descubre algún fallo lógico en el protocolo, muestra paso a paso cómo un atacante podría aprovecharlo.
